% Generated draft: requires mathematical and source review.
\documentclass[11pt]{article}
\usepackage[margin=25mm]{geometry}
\usepackage{fontspec}
\setmainfont{Latin Modern Roman}
\usepackage{amsmath,amssymb,mathtools,unicode-math}
\setmathfont{Latin Modern Math}
\usepackage{enumitem,needspace,xurl,hyperref}
\setlength{\parindent}{0pt}
\setlength{\parskip}{5pt}
\setlength{\emergencystretch}{3em}
\newcommand{\N}{\mathbb{N}}
\newcommand{\Z}{\mathbb{Z}}
\newcommand{\Q}{\mathbb{Q}}
\newcommand{\R}{\mathbb{R}}
\newcommand{\C}{\mathbb{C}}
\newcommand{\F}{\mathbb{F}}
\newcommand{\A}{\mathbb{A}}
\newcommand{\PP}{\mathbb{P}}
\newcommand{\E}{\mathbb{E}}
\newcommand{\eps}{\varepsilon}
\newcommand{\dd}{\,\mathrm{d}}
\begin{document}
\section*{UnsolvedMath ID: 11000239}
\textbf{Source problem code:} catalogue entry
\subsection{Definitions and mathematical statement}
Let $G<\operatorname{MCG}(S)$ be a finite-coarea Veech subgroup. In the customary interpretation, $G$ is induced by affine symmetries of a translation/half-translation surface, preserving a Teichm\"uller disk, with derivative Fuchsian group having finite hyperbolic covolume. Any derivative kernel or projectivization convention must agree with the source.

\textbf{Problem 5.4.} Describe the group denoted $\Gamma_C$ in the source for such a $G$. In particular, does $\Gamma_C$ contain the associated group $\Gamma_G$ as a subgroup of finite index?
\[
 \Gamma_G\le\Gamma_C,\qquad [\Gamma_C:\Gamma_G]<\infty.
\]
This is the stated index direction: the larger group is $\Gamma_C$. The excerpt does not define $C$, $\Gamma_C$, $\Gamma_G$, or the intended inclusion, so the display is a precise rendering of its target conditional on those missing constructions, not an independent definition of them.

The surrounding chapter concerns surface-group extensions. In the standard closed-surface formulation, for $\Pi=\pi_1(S)$ and a mapping-class subgroup $G\le\operatorname{Out}(\Pi)$, its pullback in $\operatorname{Aut}(\Pi)$ gives
\[
 1\longrightarrow\Pi\longrightarrow\Gamma_G\longrightarrow G\longrightarrow1.
\]
The adjacent local item 11000240 relates $\Gamma_C$ to fiber-respecting quasi-isometries of $\Gamma_G$, meaning quasi-isometries that coarsely preserve the family of cosets of the surface-group kernel. This identifies the intended coarse-geometric context, but does not supply the missing definition of $C$ or establish which precise identification applies in the present Veech case. It would be unjustified to choose an arbitrary subgroup called $C$, replace it by a Fuchsian commensurator, or infer the requested finite-index statement just from finite coarea.
\subsection{Short English statement}
For a finite-coarea Veech subgroup, determine the source's associated coarse-geometric group $\Gamma_C$, especially whether the surface-group extension $\Gamma_G$ has finite index in it. The construction of $C$ is missing from the excerpt.
\subsection{Sources}
\begin{itemize}
\item Local UnsolvedMath registry, permanent ID 11000239, source code catalogue entry; no proposer or proposal year supplied.
\item Farb, edited volume \emph{Problems on Mapping Class Groups and Related Topics}, chapter context ``Problems in the geometry of surface group extensions,'' Problem 5.4, PDF page 275; \url{https://www.math.uchicago.edu/~farb/papers/mcgbook.pdf}. The registry records PDF-font-structure extraction and access on 2026-07-29; the inventory distinguishes 329 labeled headings by source page.
\item Adjacent local ID 11000240 supplies the fiber-respecting quasi-isometry context and source label [Mos03a], but does not reproduce the definition of $C$. Full bibliography and theorem statements are absent from the present record.
\end{itemize}
\end{document}
